%%%
% UNI Beamer Template -- package imports
% Author: Yihong Liu (https://liu-yihong.github.io/)
% License: GNU GPL v3.0
%
% ENGINE: compile with LuaLaTeX (fontspec + emoji require it).
%   latexmk picks LuaLaTeX automatically via the bundled latexmkrc.
%   On Overleaf: Menu > Settings > Compiler > LuaLaTeX.
%
% NOTE: xcolor options are passed through \documentclass[xcolor={...}]
%       in the main file. Do NOT load xcolor/color/hyperref again here:
%       beamer already loads them, and reloading causes option clashes.
%%%

% ---------------------------------------------------------------- fonts
\usepackage{fontspec}
% Body font. Choose one:  fira | heros | dejavu | lmsans
%   fira   : Fira Sans -- screen-oriented, large x-height, clear I/l/1 (default)
%   heros  : TeX Gyre Heros -- free Helvetica clone
%   dejavu : DejaVu Sans -- the Bitstream Vera family that the "arev" package is based on
%   lmsans : Latin Modern Sans -- LaTeX's default; thin strokes, small x-height
% Any font that is not installed falls back silently to Latin Modern Sans.
\providecommand{\unibodyfont}{fira}

\ifdefstring{\unibodyfont}{fira}{%
  \IfFontExistsTF{FiraSans-Regular.otf}{%
    \setsansfont{FiraSans}[
      Extension = .otf, UprightFont = *-Regular, BoldFont = *-SemiBold,
      ItalicFont = *-Italic, BoldItalicFont = *-SemiBoldItalic, Scale = 0.95]%
  }{}}{}
\ifdefstring{\unibodyfont}{heros}{%
  \IfFontExistsTF{texgyreheros-regular.otf}{%
    \setsansfont{texgyreheros}[
      Extension = .otf, UprightFont = *-regular, BoldFont = *-bold,
      ItalicFont = *-italic, BoldItalicFont = *-bolditalic, Scale = 0.92]%
  }{}}{}
\ifdefstring{\unibodyfont}{dejavu}{%
  \IfFontExistsTF{DejaVuSans.ttf}{%
    \setsansfont{DejaVuSans}[
      Extension = .ttf, UprightFont = *, BoldFont = *-Bold,
      ItalicFont = *-Oblique, BoldItalicFont = *-BoldOblique, Scale = 0.86]%
  }{}}{}
\usefonttheme{professionalfonts} % keep math in the standard math font

% ---------------------------------------------------------------- core (already loaded)
% beamer and the LaTeX kernel already load these; listing them is harmless and keeps
% the full package list in one place. Pass xcolor OPTIONS via \documentclass[xcolor={...}]
% in the main file -- adding options here causes an "Option clash for package xcolor" error.
\usepackage{xcolor}
\usepackage{color}
\usepackage{hyperref}
% \usepackage[bookmarks=true,bookmarksopen=true,pdfborder={0 0 0},pdfhighlight={/N},linkbordercolor=[rgb]{0.5,0.5,0.5},implicit=false,pagebackref=true]{hyperref}
%   (hyperref is loaded by beamer; set such options with \hypersetup{...} instead)
\usepackage{textcomp}
% \usepackage[utf8]{inputenc}   % not needed with LuaLaTeX (UTF-8 is native)

% ---------------------------------------------------------------- icons & emoji
\usepackage{fontawesome5}
% \usepackage{fontawesome}      % older FontAwesome 4 icon set
\usepackage{emoji} % https://ctan.org/pkg/emoji
% Use the first color-emoji font found: Noto Color Emoji (ships with TeX Live,
% MiKTeX and Overleaf), Segoe UI Emoji (Windows), or a font file you drop into
% includes/fonts/ (e.g., joypixels-android.ttf).
\IfFontExistsTF{Noto Color Emoji}{%
  \setemojifont{Noto Color Emoji}[BoldFont=*]%
}{%
  \IfFontExistsTF{Segoe UI Emoji}{%
    \setemojifont{Segoe UI Emoji}[BoldFont=*]%
  }{%
    \IfFileExists{includes/fonts/joypixels-android.ttf}{%
      \setemojifont{joypixels-android.ttf}[Path=includes/fonts/,BoldFont=*]%
    }{}%
  }%
}

% ---------------------------------------------------------------- layout & graphics
\usepackage{setspace}
\usepackage{truncate}            % footer: shorten over-long titles with "..."
\usepackage{graphicx}
\usepackage{tikz}
\usetikzlibrary{calc}
\usepackage[justification=centering]{caption}

% ---------------------------------------------------------------- tables
\usepackage{booktabs}
\usepackage{array}
\usepackage{longtable}
\usepackage{multirow}
\usepackage{multicol}

% ---------------------------------------------------------------- text tools
\usepackage{listings}
\usepackage{soul}               % \hl{} highlighting
\usepackage[normalem]{ulem}     % \uline, \sout (keeps \emph italic)
\useunder{\uline}{\ul}{}
\usepackage{censor}             % \censor{} for answer reveals
% \usepackage{tipa}             % IPA phonetic symbols

% ---------------------------------------------------------------- bibliography
\usepackage[
  backend      = biber,
  natbib       = true,
  style        = authoryear,
  sorting      = nty,
  uniquename   = mininit,
  maxbibnames  = 99,
  maxcitenames = 2,
  mincitenames = 1,
  backref      = true,
]{biblatex}
\addbibresource{includes/bibliography/ref.bib}
